\section*{الفصل \ref{ch:g6:decomposers-and-soil} --- المحلِّلون والتربة}

\begin{solution}{exo:g6:decomposers-and-soil:1}
حبيبات معدنية (صخر مطحون)، وبقايا عضوية في كل أطوار التعفن، وماء،
وهواء --- مع جمهرتها الحية. والطبقة الداكنة هي \omterm{def:g6:decomposers-and-soil:soil}{الدُّبال}.
\end{solution}

\begin{solution}{exo:g6:decomposers-and-soil:2}
على \omterm{def:g6:plants-are-producers:matter}{المادة العضوية} الميتة --- \omterm{def:g1:plants-around-us:parts}{الأوراق} الساقطة، والخشب الميت،
والفضلات، والجثث. والمرئيون: ديدان الأرض، وحشرات القمل الخشبي،
وأمّات أربعة وأربعين، والذنبيات القافزة (أو خنافس الروث، \omterm{ex:g2:animals-grow-up:butterfly}{ويرقات}
الذباب). والمتغذي بلبّادة الخيوط: \omterm{def:g6:decomposers-and-soil:decomposers}{الفطر}.
\end{solution}

\begin{solution}{exo:g6:decomposers-and-soil:3}
عائدةً إلى \omterm{def:g6:plants-are-producers:matter}{مادة معدنية} --- دُبالًا، ثم أملاحًا معدنية مذابة، وغازًا
إلى الهواء. وهو يعكس تجارة \omterm{def:g5:food-webs:roles}{المنتِجين} التي بنت العضوي من المعدني.
\end{solution}

\begin{solution}{exo:g6:decomposers-and-soil:4}
تشرين الثاني: ساقطة كاملة. والشتاء: مليّنة منقوبة بالفطور، مخرّمة
حتى صارت هيكلًا بالذنبيات القافزة وحشرات القمل الخشبي. والربيع:
مجرورة إلى أسفل ومأكولة بدودة أرض، ومقذوفة فتاتًا. والصيف: دُبال
أسود؛ ومعادنها في ماء \omterm{def:g6:decomposers-and-soil:soil}{التربة}، تشربها \omterm{def:g1:plants-around-us:parts}{الجذور} --- ولعله \omterm{def:g1:plants-around-us:parts}{جذر} البلوطة
نفسها.
\end{solution}

\begin{solution}{exo:g6:decomposers-and-soil:5}
لأن المحلِّلين يأكلونه بالسرعة التي تُسقطه بها السنون: فكل فرش
يُفكَّك إلى دُبال ومعادن قبل أن يسقط التالي كاملًا.
\end{solution}

\begin{solution}{exo:g6:decomposers-and-soil:6}
الدفء والندى والتهوية. ويكاد العمل يتوقف في الشتاء --- من البرد
--- ولهذا يبقى الفرش كثيفًا حتى ينظّفه الربيع الدافئ الندي.
\end{solution}

\begin{solution}{exo:g6:decomposers-and-soil:7}
\omterm{def:g5:food-webs:roles}{المنتِجون} $\leftarrow$ \omterm{def:g5:food-webs:roles}{المستهلِكون} (\omterm{def:g1:plants-around-us:plant}{النباتات} تُؤكل)؛ \omterm{def:g5:food-webs:roles}{والمنتِجون}
\omterm{def:g5:food-webs:roles}{والمستهلِكون} $\leftarrow$ بقايا ميتة وفضلات \omterm{def:g1:plants-around-us:parts}{وأوراق} ساقطة
$\leftarrow$ \omterm{def:g6:decomposers-and-soil:decomposers}{المحلِّلون} (الموتى يُؤكلون)؛ \omterm{def:g6:decomposers-and-soil:decomposers}{والمحلِّلون}
$\leftarrow$ \omterm{prop:g4:what-plants-need:needs}{أملاح معدنية} في ماء \omterm{def:g6:decomposers-and-soil:soil}{التربة}
$\leftarrow$ تشربها \omterm{def:g5:food-webs:roles}{المنتِجون}: عجلة مغلقة.
\end{solution}

\begin{solution}{exo:g6:decomposers-and-soil:8}
العمّال: الديدان، وحشرات القمل الخشبي، والذنبيات القافزة،
والفطور، والأطقم غير المرئية. والمادة الخام: نفايات المطبخ
العضوية. \omterm{def:g5:food-webs:roles}{والمنتج}: دُبال داكن تغذّي معادنه الأحواض. وهي تعمل دافئة
لأن العمّال، ككل مستهلِك، يحرقون جزءًا مما يأكلون --- فالكومة
يدفّئها عيش طاقمها نفسه.
\end{solution}

\begin{solution}{exo:g6:decomposers-and-soil:9}
محلِّلة --- تجرّ \omterm{def:g1:plants-around-us:parts}{الأوراق} تحت الأرض وتأكل المادة المتعفنة ---
وحرّاثة: فأنفاقها تهوّي \omterm{def:g6:decomposers-and-soil:soil}{التربة}، وقذفاتها تعجن المعدني والعضوي
معًا. \omterm{def:g6:decomposers-and-soil:soil}{وتربة} حقول بأكملها تمرّ في أمعاء الديدان، ولهذا توّجها علماء
الطبيعة القدماء صانعةَ الأرض الزراعية.
\end{solution}

\begin{solution}{exo:g6:decomposers-and-soil:10}
الحوض المغطى: دُبال داكن هشّ؛ وإحصاء غني --- ديدان، وحشرات قمل
خشبي، وذنبيات قافزة؛ وفرش السطح يختفي بوضوح موسمًا بعد موسم؛
والحكم: دورة دائرة تستحق الإطعام. والدرب المدوس: غبار شاحب مرصوص؛
وصينية تكاد تكون خالية؛ وما يسقط يبقى بلا تعفن؛ والحكم: \omterm{def:g6:decomposers-and-soil:soil}{تربة}
جائعة بلا هواء وقد ذهب طاقمها.
\end{solution}

\begin{solution}{exo:g6:decomposers-and-soil:11}
الهواء. فالكيس المختوم المبتلّ بلا هواء يوقف العمّال المحتاجين
إليه ولا يدع يجري إلا تعفن حامض بلا هواء. والإصلاح: إخراجه من
الكيس إلى كومة سائبة رطبة مقلَّبة --- تُحفظ فيها \omterm{def:g6:exploring-our-environment:conditions}{الرطوبة} ويُدخل
الهواء.
\end{solution}

\begin{solution}{exo:g6:decomposers-and-soil:12}
ما يتكوّم: الموتى --- فرش كل خريف، وبقايا كل جيل، وجبال من \omterm{def:g6:plants-are-producers:matter}{مادة
عضوية} غير مفكَّكة. وما ينفد: أملاح \omterm{def:g6:decomposers-and-soil:soil}{التربة} المعدنية، إذ لا يعاد
تموينها متى قُطع سهم الردّ؛ فيشرب \omterm{def:g5:food-webs:roles}{المنتِجون} آخر المعادن المذابة ثم
لا يستطيعون أن يبنوا شيئًا، مهما يكن \omterm{def:g6:exploring-our-environment:conditions}{الضوء}. والشطران هما السهم
المقطوع نفسه --- \omterm{prop:g6:decomposers-and-soil:decomposition}{التحلل} --- منظورًا إليه من طرفيه.
\end{solution}

\begin{solution}{pb:g6:decomposers-and-soil:1}
\textbf{1.} لتختلف الأكياس في شيء واحد فحسب --- الشبكة --- فيمكن
أن يُنسب أي فرق في الوزن إليها: إنه الاختبار المنصف، مدفونًا.

\textbf{2.} أتستطيع حياة \omterm{def:g6:decomposers-and-soil:soil}{التربة} غير المرئية وحدها أن تفكّك
\omterm{def:g1:plants-around-us:parts}{الأوراق} --- فالكيس الخشن يخلط العمّال المرئيين وغير المرئيين؛
والكيس الناعم يعزل غير المرئيين.

\textbf{3.} أتحكم الشروطُ العملَ: \omterm{def:g1:plants-around-us:parts}{أوراق} نفسها، وشبكة نفسها، لكن
حصًى جافًّا في مقابل ضفة ندية --- فهو يختبر شرط الندى في
\cref{prop:g6:decomposers-and-soil:speed} ميدانيًّا.

\textbf{4.} الخشن يفقد أكثر، والناعم أقل، والمختوم لا شيء
تقريبًا --- والتنبؤ مأخوذ مباشرة ممن وما يُدخله كل كيس.

\textbf{5.} الخشن: عمل الطاقم كاملًا --- أسرع تفكيك. والناعم: لم
يعمل إلا غير المرئي --- فقد حقيقي لكن أبطأ. والمختوم: لا عمّال
ولا هواء --- فلا تحلل حقيقي تقريبًا، بل فساد حامض بلا هواء.

\textbf{6.} خيوط الفطور تنقب \omterm{def:g1:plants-around-us:parts}{الورقة} وتليّنها؛ والذنبيات القافزة
وحشرات القمل الخشبي تنهشها حتى تصير مخرّمة؛ والدودة تجرّ الشظايا
إلى أسفل وتأكلها وتقذف فتاتًا معالَجًا؛ والأطقم غير المرئية تنهي كل
شظية إلى دُبال ومعادن.

\textbf{7.} على أنها قوة عاملة محلِّلة حقيقية بذاتها: فقد غادر
الوزن الكيس وكل \omterm{def:g1:animals-around-us:animal}{حيوان} مرئي ممنوع، فقد استمر التفكيك على مقياس دون
البصر.

\textbf{8.} الندى قبل كل شيء --- فالحصى ينزّ ويجفّ --- ومعه
الجمهرة العاملة التي تحملها ضفة رطبة؛ \omterm{def:g6:exploring-our-environment:conditions}{فالشروط} الجافة تعطّل أي
طاقم يصل.

\textbf{9.} إلى \omterm{def:g6:decomposers-and-soil:soil}{التربة} وإلى الهواء: فجزء صار دُبالًا وأملاحًا
معدنية مذابة في ماء \omterm{def:g6:decomposers-and-soil:soil}{تربة} الضفة، وجزء أحرقه \omterm{def:g6:decomposers-and-soil:decomposers}{المحلِّلون} العاملون
فغادر غازًا --- فالغرامات المفقودة كلها محسوبة، ولم يزل منها شيء.

\textbf{10.} ليست مصادفة: فمعادن \omterm{def:g1:plants-around-us:parts}{الأوراق} المدفونة تحررت هناك
بالضبط، والقرّاص مشهور بنهمه للأرض الغنية. والشريط الأخضر هو سهم
الدورة الأخير مرئيًّا --- \omterm{prop:g6:decomposers-and-soil:decomposition}{التحلل} يغذّي \omterm{def:g5:food-webs:roles}{المنتِجين}.

\textbf{11.} لفافة المطمر هي الكيس المختوم على مقياس مدينة: فهي
بلا هواء وبلا طاقم، تمنع التحللَ عمّالَه وشروطَه، فتبقى النفايات
كما بقيت حمأة الكيس --- بينما تعود القشور نفسها ترابًا في شهور في
كومة مهوّاة ندية ملأى بالعمّال.

\textbf{12.} برهن الشتاء على أن الفرش ينظّفه محلِّلو \omterm{def:g6:decomposers-and-soil:soil}{التربة} ---
الطاقم المرئي والأطقم غير المرئية معًا، وأسرع ما يكون في الدفء
والندى والتهوية. وتنظيفهم يغذّي مخازن \omterm{def:g6:decomposers-and-soil:soil}{التربة} المعدنية، ومن خلالها
\omterm{def:g5:food-webs:roles}{المنتِجين}: إنها عجلة المادة، تدور في ضفة سياج.
\end{solution}
